% !TeX root = main.tex
\documentclass[12pt, a4paper]{project} % Dùng class tùy chỉnh có định nghĩa \coverpage và các biến trang bìa

% ==========================================
% 1. CÀI ĐẶT NGÔN NGỮ & FONT CHỮ
% ==========================================
\usepackage[T5]{fontenc}      % Hỗ trợ font tiếng Việt
\usepackage[utf8]{inputenc}   % Hỗ trợ gõ Unicode UTF-8
\usepackage[vietnamese]{babel} % Định dạng ngày tháng, tự động dịch các tiêu đề (Mục lục, Chương...)
\AtBeginDocument{\selectlanguage{vietnamese}}
\addto\captionsvietnamese{\renewcommand{\chaptername}{Chương}}

% Tăng cỡ chữ cho tiêu đề chương (vd: "Giới thiệu")
\makechapterstyle{nt531chapter}{
    \chapterstyle{default}
    \renewcommand{\chapnamefont}{\normalfont\Large\bfseries}
    \renewcommand{\chapnumfont}{\normalfont\Huge\bfseries}
    \renewcommand{\chaptitlefont}{\normalfont\Huge\bfseries}
}
\chapterstyle{nt531chapter}

% ==========================================
% 2. CĂN LỀ & ĐỊNH DẠNG TRANG
% ==========================================
\geometry{top=2.5cm, bottom=2.5cm, left=3cm, right=2cm} % Class đã load geometry, chỉ set lại thông số
\usepackage{fancyhdr}         % Tùy biến header/footer (đánh số trang)
\setlength{\parskip}{0.5em}   % Khoảng cách giữa các đoạn văn
\linespread{1.3}              % Dãn dòng 1.3 (tương đương 1.5 lines trong Word)

% ==========================================
% 3. XỬ LÝ BẢNG BIỂU (TABLES)
% ==========================================
\usepackage{tabularx}         % Tự động dàn trang bảng cho vừa chiều rộng
\usepackage{booktabs}         % Đường kẻ bảng chuyên nghiệp (toprule, midrule, bottomrule)
\usepackage{multirow}         % Gộp dòng (hữu ích cho bảng RQ/Fairness)
\usepackage{makecell}         % Cho phép ngắt dòng (enter) ngay trong 1 ô của bảng
\usepackage{longtable}        % Cho phép bảng tự động ngắt trang nếu quá dài (hữu ích cho Workloads table)

% ==========================================
% 4. XỬ LÝ HÌNH ẢNH & BIỂU ĐỒ (FIGURES)
% ==========================================
\usepackage{graphicx}         % Gói cơ bản để chèn ảnh
\usepackage{float}            % Ép vị trí ảnh bằng lệnh [H]
\usepackage{subcaption}       % Cho phép chèn nhiều ảnh nhỏ trong 1 ảnh lớn (VD: Đồ thị so sánh p99)

% ==========================================
% 5. CÔNG THỨC TOÁN HỌC & THỐNG KÊ
% ==========================================
\usepackage{amsmath, amssymb, amsfonts} % Hỗ trợ gõ công thức (p50, p99, O(1), O(n))

% ==========================================
% 6. MÃ NGUỒN & LOG (CODE & ARTIFACTS)
% ==========================================
\usepackage{listings}         % Chèn code (hữu ích cho manifests, bash script, hubble.log)
\usepackage{xcolor}           % Hỗ trợ tô màu cho code snippet

% Cấu hình mặc định cho khối code
\lstset{
    basicstyle=\ttfamily\small,
    breaklines=true,
    frame=single,
    backgroundcolor=\color{gray!10},
    keywordstyle=\color{blue},
    commentstyle=\color{green!50!black},
    stringstyle=\color{red}
}

% ==========================================
% 7. DANH SÁCH & LIÊN KẾT (LISTS & HYPERLINKS)
% ==========================================
\usepackage{enumitem}         % Tùy biến bullet points và numbering
\usepackage{hyperref}         % Class đã load hyperref; load lại không thêm options để tránh clash
\hypersetup{hidelinks}        % Ẩn viền link trong Mục lục/Danh mục
\usepackage{cleveref}         % Tham chiếu thông minh (tự động thêm chữ "Bảng", "Hình" trước số)
